\subsection{Robust policy search with bounds}\label{\refsec:RobustRoot}
The politico-economic equilibria below are identified from first order conditions rather than by maximising political objectives directly. This is paramount to keeping the state space small --- it is what allows the closed-form reduction of the predetermined savings shares to do its work --- but it converts a maximisation problem into a root problem, and a root problem inherits none of the safeguards of the maximisation it came from: a root may be a minimum, there may be several, and the optimum may sit at a corner where no root exists at all. This subsection collects, once, the machinery that restores those safeguards. It is stated for a generic first order condition and is implemented as a standalone package (\texttt{python/gridsearch} in the repository); the model enters only through the function being evaluated.

Throughout, let $f(\tau)$ denote a first order condition --- in our applications the derivative of a political objective $\mathcal{W}$ with respect to the tax rate --- that it is safe to evaluate on an interval $[l,u]$ with $0\leq l < u < 1$, and let $z(\tau) = f\left([\tau]_{[l,u]}\right)$ denote its bounded evaluation, where $[\tau]_{[l,u]} \equiv \min(u, \max(l,\tau))$.

\smalltitle{A bounded-root reparameterization.}
Gradient-based solvers are efficient when they apply, but they search all of $\mathbb{R}$ and $f$ is only meaningful on $[l,u]$. Instead of solving $f(\tau) = 0$ subject to bounds, one can solve the unconstrained auxiliary problem
\begin{align}\label{\refeq:root}
    h(z,\tau) = z - |z| \underbrace{\left[a_0\min\left(\tau - l,0\right) + a_1\max\left(\tau - u, 0\right)\right]}_{\equiv g(\tau)} = 0,
\end{align}
where $a_0, a_1>0$ are technical parameters. This gives three regions:
\begin{align*}
    h\left(z,\tau\right) = \left\lbrace \begin{array}{ll} f(l)-|f(l)| a_0 (\tau - l), & \tau<l \\
    f(\tau), & \tau\in[l,u] \\
    f(u)-|f(u)|a_1(\tau-u), & \tau>u\end{array}\right., && \dfrac{\partial h}{\partial \tau} = \left\lbrace \begin{array}{ll} -|f(l)|a_0, & \tau<l \\
    \dfrac{\partial f}{\partial \tau}, & \tau\in[l,u] \\
    -|f(u)|a_1, & \tau>u\end{array}\right.
\end{align*}
$f$ is never evaluated outside $[l,u]$, yet $h$ is defined everywhere --- affine with a weakly negative slope outside the interval --- so an unconstrained solver can search freely and the bounded solution is recovered as $\tau = [\tilde{\tau}]_{[l,u]}$. The construction also detects corners: if $f>0$ throughout $[l,u]$, the auxiliary problem still has a root, at $\tilde{\tau} = u+1/a_1$, which bounds to the corner solution $\tau = u$; symmetrically $f<0$ throughout yields $\tilde{\tau} = l-1/a_0$ and thus $\tau = l$. For completeness we also record the \emph{extended grid} construction that makes these corners visible to a grid-based root detector: given an interior grid (below), add two outer nodes
\begin{align}\label{\refeq:extendedGrid}
    \tilde{\tau}^{(0)} = l-\frac{1}{a_0}-\delta, \qquad \tilde{\tau}^{(M+1)} = u+\frac{1}{a_1}+\delta,
\end{align}
for a small $\delta>0$, sitting a hair beyond the two boundary roots of \eqref{\refeq:root}. Evaluating $h$ at $\tilde{\tau} = l-1/a_0$ gives $h = f(l)+|f(l)|$, identically zero whenever $f(l)<0$, and symmetrically at the upper node: a corner is encoded as an \emph{exact} zero rather than as a sign change, so any root detection must treat exact zeros as first-class citizens rather than test only for $h^{(k)}h^{(k+1)}<0$. In this model the grid searches themselves do not need the extended grid --- the selection rule below places $l$ and $u$ in the candidate set directly, which is precisely what the outer nodes exist to achieve --- so \eqref{\refeq:root} and \eqref{\refeq:extendedGrid} serve here as the bounded-evaluation convention and as an optional gradient-based polishing device, not as part of the grid construction.

\smalltitle{The grid, and selection among several roots.}
Let $\mathcal{T} = \left\lbrace \tau^{(1)},...,\tau^{(M)}\right\rbrace$ with $\tau^{(1)} = l$ and $\tau^{(M)} = u$ denote a grid on the interior, on which $z$ is evaluated everywhere. Solving $z = 0$ is only a necessary condition for the political optimum: an upward crossing of $z$ is a local \emph{minimum} of $\mathcal{W}$, and when several downward crossings exist the first one located need not be the global maximum. Since the grid search delivers $z$ everywhere on $\mathcal{T}$ in any case, the maximisation can be restored without any additional evaluations of $z$.

Let $\hat{z}$ denote the piecewise-linear interpolant of $z$ through the nodes of $\mathcal{T}$ --- the same interpolant that is used to locate crossings --- and define
\begin{align}\label{\refeq:objectiveProfile}
    \hat{\mathcal{V}}(\tau) = \int_l^{\tau}\hat{z}(x)\mbox{ }dx.
\end{align}
Because $z = d\mathcal{W}/d\tau$, $\hat{\mathcal{V}}$ approximates the political objective up to an additive constant. It is piecewise quadratic and continuously differentiable with $\hat{\mathcal{V}}' = \hat{z}$, so its interior local maxima are exactly the downward crossings of $\hat{z}$, and
\begin{align}\label{\refeq:candidates}
    \arg\max_{\tau\in[l,u]}\hat{\mathcal{V}}(\tau)\in\mathcal{C} \equiv \left\lbrace l,u \right\rbrace\cup\left\lbrace \tau\in(l,u)\mbox{ }:\mbox{ }\hat{z}(\tau) = 0\text{ and }\hat{z}\text{ crosses from above}\right\rbrace.
\end{align}
The candidate set $\mathcal{C}$ is finite and cheap to construct, and $\hat{\mathcal{V}}$ is evaluated on it \emph{exactly}: at the nodes by the trapezoidal rule,
\begin{align*}
    \hat{\mathcal{V}}\left(\tau^{(k+1)}\right) = \hat{\mathcal{V}}\left(\tau^{(k)}\right)+\frac{1}{2}\left(z\left(\tau^{(k)}\right)+z\left(\tau^{(k+1)}\right)\right)\left(\tau^{(k+1)}-\tau^{(k)}\right),
\end{align*}
and, for a crossing $\tau^*$ located inside the cell $\left[\tau^{(k)},\tau^{(k+1)}\right]$, by adding the triangle $\frac{1}{2}z\left(\tau^{(k)}\right)\left(\tau^*-\tau^{(k)}\right)$. Setting $\tau = \arg\max_{\tau\in\mathcal{C}}\hat{\mathcal{V}}(\tau)$ resolves multiplicity and corner solutions under a single criterion: $l$ and $u$ compete with the interior crossings on the same footing, so the rule subsumes the boundary check rather than supplementing it. The accuracy demanded of \eqref{\refeq:objectiveProfile} is only that it rank well-separated candidates correctly, which is a far weaker requirement than the accuracy demanded of the located roots themselves. When a feasibility mask excludes part of $\mathcal{T}$ (as it will below), the rule is applied to the feasible sub-grid, with $l$ and $u$ replaced by its endpoints.

\smalltitle{Root problems that are not maximisations.}
Not every root problem below is a first order condition. The state fixed points of \cref{\refsec:PEELOG} and \cref{\refsec:PEE} are self-consistency conditions in which \emph{any} crossing is a solution; for those, the direction filter and the objective integration above are deliberately not applied. Keeping the two kinds of problem distinct matters, because the safeguards appropriate to one are errors in the other.

\smalltitle{Two implementation conventions.}
First, convergence of any gradient-polished solution is judged on $\max\left|z\right|$ directly, against an explicit tolerance, rather than on the solver's own success flag --- the flags reported by standard library routines are not consistent across methods, and a residual check is both stricter and portable. Second, a residual is only a meaningful convergence measure at an \emph{interior} optimum: at a corner, $z\neq 0$ is the correct answer, not an error, so residual diagnostics are restricted to periods (or states) whose selected policy is interior.
